% 例 1：作点 A 关于直线 l 的对称点 A'
\documentclass[tikz,border=4pt]{standalone}
\usepackage{ctex}
\usepackage{amsmath}
\usetikzlibrary{calc, decorations.markings}
\begin{document}
\begin{tikzpicture}[scale=1.0, thick,
  tickone/.style={postaction={decorate, decoration={markings,
    mark=at position 0.5 with {\draw (-1.5pt,-3pt) -- (1.5pt,3pt);}}}}]
  \draw (-3,0) -- (3,0) node[right]{$l$};
  \coordinate (A) at (-1.5,1.8);
  \coordinate (M) at (-1.5,0);
  \coordinate (Ap) at (-1.5,-1.8);
  \draw[red, thick] (A) -- (M);
  \draw[red, thick, dashed] (M) -- (Ap);
  \draw[tickone] (A) -- (M);
  \draw[tickone] (M) -- (Ap);
  % 直角标记
  \draw (-1.7,0) -- (-1.7,0.2) -- (-1.5,0.2);
  \fill (A) circle (1.2pt); \fill (Ap) circle (1.2pt); \fill (M) circle (1.0pt);
  \node[above] at (A) {$A$};
  \node[below] at (Ap) {$A'$};
  \node[below right] at (M) {\footnotesize $M$};
\end{tikzpicture}
\end{document}
